\scanpage{35}\mspage{34}

\uline{20 juin} Comme
\[
\mathrm{Fig}(M)\longrightarrow\mathrm{Fig}(L)
\]
induit par l'application
\[
L\longrightarrow M
\]
est croissante pour $\leq$ et $\ll$, il s'ensuit que l'application
\[
\begin{aligned}
{\unsure{\mathcal F|\mathscr F}}&\longrightarrow\mathrm{Fig}(L) & F&\longrightarrow\unsure{\mathrm{Mulomb}|\mathrm{mulomb}}(F)\\
&&&(\uline{\text{petite}}\ \text{multiombre})
\end{aligned}
\]
est également croissante pour $\leq$ et $\ll$, et \unsure{commute} aux $\mathrm{Sup}$ pour $\leq$, et à $\mathrm{Inf}$ pour $\leq$. Mais elle n'est plus \unsure{néc|néc.} injective.
\margin{[\unsure{transforme} sous-figure en sous-figure.]}

\uline{Exemple 1}\quad \deleted{\unsure{a|1}}A) Soit $\deleted{X}E$ espace topologique, $\deleted{\unsure{A|M}}$
$\mathcal F$ l'ensemble des \deleted{figures ensemblistes} \deleted{paires} triples
\[
F=(X,\deleted{\unsure{\sigma}}\Sigma,\unsure{F|\Phi})
\]
où
\[
\begin{aligned}
&X\ \text{est une partie fermée de } E\\
&\Sigma\ \text{est une structure \unsure{lin.|liss.} par morceaux (disons, ou}\\
&\qquad\text{``\unsure{modérée}'' pour une théorie donnée) sur } X\\
&\Phi\ \text{est une figure ens. sur } X,\ \text{de support }X,\\
&\qquad\text{dont les strates sont des sous-ens. fermés}\\
&\qquad\text{lin. par morceaux.}
\end{aligned}
\]
\[
\begin{aligned}
(X',\Sigma',\Phi') &\leq (X^{\deleted{\unsure{\prime}}},\Sigma^{\deleted{\unsure{\prime}}},\Phi^{\deleted{\unsure{\prime}}})
\overset{\mathrm{dfn}}{\Longleftrightarrow}
\left\{\begin{aligned}
&\text{a) }X'\subset X\unsure{,|{}_{s}},\ \text{sous-espace lin.}\\
&\hspace{4em}\text{par morceaux pour }\Sigma\\
&\text{b) }\deleted{\unsure{\Sigma'} }\;\Sigma'\ \text{induite par }\Sigma\\
&\text{c) }\Phi'\leq\Phi
\end{aligned}\right.\\
(X',\Sigma',\Phi') &\ll (X,\Sigma,\Phi)
\overset{\mathrm{dfn}}{\Longleftrightarrow}
\left\{\begin{aligned}
&\text{a) et b) comme dessus}\\
&\text{c) }\Phi'\ll\Phi
\end{aligned}\right.
\end{aligned}
\]
On vérifie les axiomes. \unsure{L'ens.|D'ail.} $L$ s'identifie à $\deleted{X}E$, et l'application \unsure{en} $\mathcal F\longrightarrow\mathrm{Fig}(\deleted{\unsure{E}}L)=\mathrm{Fig}(E)$ est l'application $F\deleted{\mapsto}=(X,\Sigma,\Phi)\longmapsto\Phi$. Elle n'est visiblement pas injective en \unsure{générale|général|généralité} (p. ex. si $\deleted{\unsure{X}}E$ contient partie hom. à $[0,1]$).

B) \uline{Variante}\quad $M={}$ ens des \deleted{\unsure{parties de}} \inserted{couples} $X=(\deleted{X}A,\sigma)$ avec :
\[
\left\{\begin{aligned}
&\deleted{X}A\ \text{partie \deleted{\unsure{finie fermée|fermée|finie}} de }\deleted{X}E\\
&\sigma\ \text{une structure de \uline{simplexe affine} sur }\deleted{X}A\\
&\quad\text{d'où figure ensembliste }\Phi_\sigma\ \text{\unsure{sur|dans} }|\Phi_\sigma|=\deleted{X}A\\
&\quad\text{\unsure{son|de} support. l'inclusion }\deleted{X}A\subset E\ \text{continue.}
\end{aligned}\right.
\]
